%%%%%%%%%%%%%%%%%%%%%%%%%%%%%%%%%%%%%%%%%%%%%%%%%%%%%%%%%%%%%%%%%%%%%%
% Template for a UBC-compliant dissertation
% At the minimum, you will need to change the information found
% after the "Document meta-data"
%
%!TEX TS-program = pdflatex
%!TEX encoding = UTF-8 Unicode

%% The ubcdiss class provides several options:
%%   gpscopy (aka fogscopy)
%%       set parameters to exactly how GPS specifies
%%         * single-sided
%%         * page-numbering starts from title page
%%         * the lists of figures and tables have each entry prefixed
%%           with 'Figure' or 'Table'
%%       This can be tested by `\ifgpscopy ... \else ... \fi'
%%   10pt, 11pt, 12pt
%%       set default font size
%%   oneside, twoside
%%       whether to format for single-sided or double-sided printing
%%   balanced
%%       when double-sided, ensure page content is centred
%%       rather than slightly offset (the default)
%%   singlespacing, onehalfspacing, doublespacing
%%       set default inter-line text spacing; the ubcdiss class
%%       provides \textspacing to revert to this configured spacing
%%   draft
%%       disable more intensive processing, such as including
%%       graphics, etc.
%%

% For submission to GPS
\documentclass[gpscopy,onehalfspacing,11pt]{ubcdiss}

% For your own copies (looks nicer)
% \documentclass[balanced,twoside,11pt]{ubcdiss}

%%%%%%%%%%%%%%%%%%%%%%%%%%%%%%%%%%%%%%%%%%%%%%%%%%%%%%%%%%%%%%%%%%%%%%
%%%%%%%%%%%%%%%%%%%%%%%%%%%%%%%%%%%%%%%%%%%%%%%%%%%%%%%%%%%%%%%%%%%%%%
%%
%% FONTS:
%% 
%% The defaults below configures Times Roman for the serif font,
%% Helvetica for the sans serif font, and Courier for the
%% typewriter-style font.  Configuring fonts can be time
%% consuming; we recommend skipping to END FONTS!
%% 
%% If you're feeling brave, have lots of time, and wish to use one
%% your platform's native fonts, see the commented out bits below for
%% XeTeX/XeLaTeX.  This is not for the faint at heart. 
%% (And shouldn't you be writing? :-)
%%

%% NFSS font specification (New Font Selection Scheme)
\usepackage{times,mathptmx,courier}
\usepackage[scaled=.92]{helvet}

%% Math or theory people may want to include the handy AMS macros
%\usepackage{amssymb}
%\usepackage{amsmath}
%\usepackage{amsfonts}

%% The pifont package provides access to the elements in the dingbat font.   
%% Use \ding{##} for a particular dingbat (see p7 of psnfss2e.pdf)
%%   Useful:
%%     51,52 different forms of a checkmark
%%     54,55,56 different forms of a cross (saltyre)
%%     172-181 are 1-10 in open circle (serif)
%%     182-191 are 1-10 black circle (serif)
%%     192-201 are 1-10 in open circle (sans serif)
%%     202-211 are 1-10 in black circle (sans serif)
%% \begin{dinglist}{##}\item... or dingautolist (which auto-increments)
%% to create a bullet list with the provided character.
\usepackage{pifont}

%%%%%%%%%%%%%%%%%%%%%%%%%%%%%%%%%%%%%%%%%%%%%%%%%%%%%%%%%%%%%%%%%%%%%%
%% Configure fonts for XeTeX / XeLaTeX using the fontspec package.
%% Be sure to check out the fontspec documentation.
%\usepackage{fontspec,xltxtra,xunicode}	% required
%\defaultfontfeatures{Mapping=tex-text}	% recommended
%% Minion Pro and Myriad Pro are shipped with some versions of
%% Adobe Reader.  Adobe representatives have commented that these
%% fonts can be used outside of Adobe Reader.
%\setromanfont[Numbers=OldStyle]{Minion Pro}
%\setsansfont[Numbers=OldStyle,Scale=MatchLowercase]{Myriad Pro}
%\setmonofont[Scale=MatchLowercase]{Andale Mono}

%% Other alternatives:
%\setromanfont[Mapping=tex-text]{Adobe Caslon}
%\setsansfont[Scale=MatchLowercase]{Gill Sans}
%\setsansfont[Scale=MatchLowercase,Mapping=tex-text]{Futura}
%\setmonofont[Scale=MatchLowercase]{Andale Mono}
%\newfontfamily{\SYM}[Scale=0.9]{Zapf Dingbats}
%% END FONTS
%%%%%%%%%%%%%%%%%%%%%%%%%%%%%%%%%%%%%%%%%%%%%%%%%%%%%%%%%%%%%%%%%%%%%%
%%%%%%%%%%%%%%%%%%%%%%%%%%%%%%%%%%%%%%%%%%%%%%%%%%%%%%%%%%%%%%%%%%%%%%



%%%%%%%%%%%%%%%%%%%%%%%%%%%%%%%%%%%%%%%%%%%%%%%%%%%%%%%%%%%%%%%%%%%%%%
%%%%%%%%%%%%%%%%%%%%%%%%%%%%%%%%%%%%%%%%%%%%%%%%%%%%%%%%%%%%%%%%%%%%%%
%%
%% Recommended packages
%%
\usepackage{checkend}	% better error messages on left-open environments
\usepackage{graphicx}	% for incorporating external images

%% booktabs: provides some special commands for typesetting tables as used
%% in excellent journals.  Ignore the examples in the Lamport book!
\usepackage{booktabs}

%% listings: useful support for including source code listings, with
%% optional special keyword formatting.  The \lstset{} causes
%% the text to be typeset in a smaller sans serif font, with
%% proportional spacing.
\usepackage{listings}
\lstset{basicstyle=\sffamily\scriptsize,showstringspaces=false,fontadjust}

%% The acronym package provides support for defining acronyms, providing
%% their expansion when first used, and building glossaries.  See the
%% example in glossary.tex and the example usage throughout the example
%% document.
%% NOTE: to use \MakeTextLowercase in the \acsfont command below,
%%   we *must* use the `nohyperlinks' option -- it causes errors with
%%   hyperref otherwise.  See Section 5.2 in the ``LaTeX 2e for Class
%%   and Package Writers Guide'' (clsguide.pdf) for details.
\usepackage[printonlyused,nohyperlinks]{acronym}
%% The ubcdiss.cls loads the `textcase' package which provides commands
%% for upper-casing and lower-casing text.  The following causes
%% the acronym package to typeset acronyms in small-caps
%% as recommended by Bringhurst.
\renewcommand{\acsfont}[1]{{\scshape \MakeTextLowercase{#1}}}

%% color: add support for expressing colour models.  Grey can be used
%% to great effect to emphasize other parts of a graphic or text.
%% For an excellent set of examples, see Tufte's "Visual Display of
%% Quantitative Information" or "Envisioning Information".
\usepackage{color}
\definecolor{greytext}{gray}{0.5}

%% comment: provides a new {comment} environment: all text inside the
%% environment is ignored.
%%   \begin{comment} ignored text ... \end{comment}
\usepackage{comment}

%% The natbib package provides more sophisticated citing commands
%% such as \citeauthor{} to provide the author names of a work,
%% \citet{} to produce an author-and-reference citation,
%% \citep{} to produce a parenthetical citation.
%% We use \citeeg{} to provide examples
\usepackage[numbers,sort&compress]{natbib}
\newcommand{\citeeg}[1]{\citep[e.g.,][]{#1}}

%% The titlesec package provides commands to vary how chapter and
%% section titles are typeset.  The following uses more compact
%% spacings above and below the title.  The titleformat that follow
%% ensure chapter/section titles are set in singlespace.
\usepackage[compact]{titlesec}
\titleformat*{\section}{\singlespacing\raggedright\bfseries\Large}
\titleformat*{\subsection}{\singlespacing\raggedright\bfseries\large}
\titleformat*{\subsubsection}{\singlespacing\raggedright\bfseries}
\titleformat*{\paragraph}{\singlespacing\raggedright\itshape}

%% The caption package provides support for varying how table and
%% figure captions are typeset.
\usepackage[format=hang,indention=-1cm,labelfont={bf},margin=1em]{caption}

%% url: for typesetting URLs and smart(er) hyphenation.
%% \url{http://...} 
\usepackage{url}
\urlstyle{sf}	% typeset urls in sans-serif


%%%%%%%%%%%%%%%%%%%%%%%%%%%%%%%%%%%%%%%%%%%%%%%%%%%%%%%%%%%%%%%%%%%%%%
%%%%%%%%%%%%%%%%%%%%%%%%%%%%%%%%%%%%%%%%%%%%%%%%%%%%%%%%%%%%%%%%%%%%%%
%%
%% Possibly useful packages: you may need to explicitly install
%% these from CTAN if they aren't part of your distribution;
%% teTeX seems to ship with a smaller base than MikTeX and MacTeX.
%%
%\usepackage{pdfpages}	% insert pages from other PDF files
%\usepackage{longtable}	% provide tables spanning multiple pages
%\usepackage{chngpage}	% support changing the page widths on demand
%\usepackage{tabularx}	% an enhanced tabular environment

%% enumitem: support pausing and resuming enumerate environments.
%\usepackage{enumitem}

%% rotating: provides two environments, sidewaystable and sidewaysfigure,
%% for typesetting tables and figures in landscape mode.  
%\usepackage{rotating}

%% subfig: provides for including subfigures within a figure,
%% and includes being able to separately reference the subfigures.
%\usepackage{subfig}

%% ragged2e: provides several new new commands \Centering, \RaggedLeft,
%% \RaggedRight and \justifying and new environments Center, FlushLeft,
%% FlushRight and justify, which set ragged text and are easily
%% configurable to allow hyphenation.
%\usepackage{ragged2e}

%% The ulem package provides a \sout{} for striking out text and
%% \xout for crossing out text.  The normalem and normalbf are
%% necessary as the package messes with the emphasis and bold fonts
%% otherwise.
%\usepackage[normalem,normalbf]{ulem}    % for \sout

%%%%%%%%%%%%%%%%%%%%%%%%%%%%%%%%%%%%%%%%%%%%%%%%%%%%%%%%%%%%%%%%%%%%%%
%% HYPERREF:
%% The hyperref package provides for embedding hyperlinks into your
%% document.  By default the table of contents, references, citations,
%% and footnotes are hyperlinked.
%%
%% Hyperref provides a very handy command for doing cross-references:
%% \autoref{}.  This is similar to \ref{} and \pageref{} except that
%% it automagically puts in the *type* of reference.  For example,
%% referencing a figure's label will put the text `Figure 3.4'.
%% And the text will be hyperlinked to the appropriate place in the
%% document.
%%
%% Generally hyperref should appear after most other packages

%% The following puts hyperlinks in very faint grey boxes.
%% The `pagebackref' causes the references in the bibliography to have
%% back-references to the citing page; `backref' puts the citing section
%% number.  See further below for other examples of using hyperref.
%% 2009/12/09: now use `linktocpage' (Jacek Kisynski): GPS now prefers
%%   that the ToC, LoF, LoT place the hyperlink on the page number,
%%   rather than the entry text.
\usepackage[bookmarks,bookmarksnumbered,%
    allbordercolors={0.8 0.8 0.8},%
    pagebackref,linktocpage%
    ]{hyperref}
%% The following change how the the back-references text is typeset in a
%% bibliography when `backref' or `pagebackref' are used
\renewcommand\backrefpagesname{\(\rightarrow\) pages}
\renewcommand\backref{\textcolor{greytext} \backrefpagesname\ }

%% The following uses most defaults, which causes hyperlinks to be
%% surrounded by colourful boxes; the colours are only visible in
%% PDFs and don't show up when printed:
%\usepackage[bookmarks,bookmarksnumbered]{hyperref}

%% The following disables the colourful boxes around hyperlinks.
%\usepackage[bookmarks,bookmarksnumbered,pdfborder={0 0 0}]{hyperref}

%% The following disables all hyperlinking, but still enabled use of
%% \autoref{}
%\usepackage[draft]{hyperref}

%% The following commands causes chapter and section references to
%% uppercase the part name.
\renewcommand{\chapterautorefname}{Chapter}
\renewcommand{\sectionautorefname}{Section}
\renewcommand{\subsectionautorefname}{Section}
\renewcommand{\subsubsectionautorefname}{Section}

%% If you have long page numbers (e.g., roman numbers in the 
%% preliminary pages for page 28 = xxviii), you might need to
%% uncomment the following and tweak the \@pnumwidth length
%% (default: 1.55em).  See the tocloft documentation at
%% http://www.ctan.org/tex-archive/macros/latex/contrib/tocloft/
% \makeatletter
% \renewcommand{\@pnumwidth}{3em}
% \makeatother

%%%%%%%%%%%%%%%%%%%%%%%%%%%%%%%%%%%%%%%%%%%%%%%%%%%%%%%%%%%%%%%%%%%%%%
%%%%%%%%%%%%%%%%%%%%%%%%%%%%%%%%%%%%%%%%%%%%%%%%%%%%%%%%%%%%%%%%%%%%%%
%%
%% Some special settings that controls how text is typeset
%%
% \raggedbottom		% pages don't have to line up nicely on the last line
% \sloppy		% be a bit more relaxed in inter-word spacing
% \clubpenalty=10000	% try harder to avoid orphans
% \widowpenalty=10000	% try harder to avoid widows
% \tolerance=1000

%% And include some of our own useful macros
\input{macros}

%%%
% ---------- JOEY'S PACKAGES ---------- %
	% Keeps all the figures from jumping into other section
	\usepackage{placeins}
	\usepackage{amsmath}
	\usepackage{tabularx}
	\newcolumntype{L}[1]{>{\raggedright\arraybackslash}p{#1}}
	\usepackage{rotating}
	\usepackage{upgreek}
	\usepackage{longtable}
	\usepackage{subcaption}
	\usepackage{pdflscape}
	\usepackage[letterpaper,left=1.25in,right=1in,bottom=1in,top=1in]{geometry}



%%%%%%%%%%%%%%%%%%%%%%%%%%%%%%%%%%%%%%%%%%%%%%%%%%%%%%%%%%%%%%%%%%%%%%
%%%%%%%%%%%%%%%%%%%%%%%%%%%%%%%%%%%%%%%%%%%%%%%%%%%%%%%%%%%%%%%%%%%%%%
%%
%% Document meta-data: be sure to also change the \hypersetup information
%%

%---------------------------------------------------------------------------%
%% ------------------------ Thesis Title --------------------- %%
\title{The title of my thesis with only all lower case except for the first word}
%\subtitle{If you want a subtitle}

% Name
\author{Jozo Texino}
\previousdegree{BA, Previous Degree, Institution, YYYY}

% Degree title
\degreetitle{Master of Science}
\institution{The University of British Columbia}
\campus{Vancouver}
\faculty{The Faculty of Graduate and Postdoctoral Studies}
\department{Geography}
\submissionmonth{August}
\submissionyear{YYYY}

%% hyperref package provides support for embedding meta-data in .PDF
%% files
\hypersetup{
pdftitle={The title of my thesis with only all lower case except for the first word (DRAFT: \today)},
pdfauthor={Jozo Texino},
pdfkeywords={key words, key words, key words, key words key words, key words, key words}
}

%%%%%%%%%%%%%%%%%%%%%%%%%%%%%%%%%%%%%%%%%%%%%%%%%%%%%%%%%%%%%%%%%%%%%%
%%%%%%%%%%%%%%%%%%%%%%%%%%%%%%%%%%%%%%%%%%%%%%%%%%%%%%%%%%%%%%%%%%%%%%
%% 
%% The document content
%%

%% LaTeX's \includeonly commands causes any uses of \include{} to only
%% include files that are in the list.  This is helpful to produce
%% subsets of your thesis (e.g., for committee members who want to see
%% the dissertation chapter by chapter).  It also saves time by 
%% avoiding reprocessing the entire file.
%\includeonly{intro,conclusions}
%\includeonly{discussion}

\begin{document}

%---------------------------------------------------------------------------%
% The document content
%---------------------------------------------------------------------------%

%% LaTeX's \includeonly commands causes any uses of \include{} to only
%% include files that are in the list.  This is helpful to produce
%% subsets of your thesis (e.g., for committee members who want to see
%% the dissertation chapter by chapter).  It also saves time by 
%% avoiding reprocessing the entire file.
%\includeonly{intro,conclusions}
%\includeonly{discussion}

	% 1. Title page (mandatory)
	\maketitle

	% 2. Abstract (mandatory - maximum 350 words)
	\chapter{Abstract}

% I think the maximum is something like 350 words
Lorem ipsum dolor sit amet, consectetur adipiscing elit, sed do eiusmod tempor incididunt ut labore et dolore magna aliqua. Ut enim ad minim veniam, quis nostrud exercitation ullamco laboris nisi ut aliquip ex ea commodo consequat. Duis aute irure dolor in reprehenderit in voluptate velit esse cillum dolore eu fugiat nulla pariatur. Excepteur sint occaecat cupidatat non proident, sunt in culpa qui officia deserunt mollit anim id est laborum.
 

	\cleardoublepage

	% 3. Preface
	\chapter{Preface}

This thesis is original work completed by the author.  Guidance was given by the supervisory committee and field assistance was provided by PERSON1, PERSON1, PERSON1, PERSON1, PERSON1, PERSON1, PERSON1, PERSON1, PERSON1, and PERSON1.

The ACTIVITY described in \autoref{} was completed in collaboration with PERSON1 (ROLE/TITLE), PERSON1, and PERSON1. 

Photos were provided by PERSON1 andPERSON1. 

A version of the work in \autoref{} and \autoref{} has been published as a poster [Last Name, First Initial., Last Name, First Initial., Last Name, First Initial. Title of the Poster]. The author acted as lead investigator, composing and presenting the poster at the CONFERENCE NAME YYYY, CITY, STATE/COUNTRY.

A version of the work in \autoref{} and \autoref{} has been published as a TALK [Last Name, First Initial., Last Name, First Initial., Last Name, First Initial. Title of the Poster]. The author acted as lead investigator, composing and presenting the TALK at the CONFERENCE NAME YYYY, CITY, STATE/COUNTRY.

	\cleardoublepage

	% 4. Table of contents (mandatory - list all items in the preliminary pages starting with the abstract, followed by chapter headings and subheadings, bibliographies and appendices)
	\tableofcontents
	\cleardoublepage	% required by tocloft package

	% 5. List of tables (mandatory if thesis has tables)
	\listoftables
	\cleardoublepage	% required by tocloft package

	% 6. List of figures (mandatory if thesis has figures)
	\listoffigures
	\cleardoublepage	% required by tocloft package

	% 7. List of illustrations (mandatory if thesis has illustrations)
	% 8. Lists of symbols, abbreviations or other (optional)

	% 9. Glossary (optional)
	%% The following is a directive for TeXShop to indicate the main file
%%!TEX root = diss.tex

\chapter{Acronyms}

% This glossary uses the handy \latexpackage{acroynym} package to automatically
% maintain the glossary.  It uses the package's \texttt{printonlyused}
% option to include only those acronyms explicitly referenced in the
% \LaTeX\ source.

% use \acrodef to define an acronym, but no listing
\acrodef{UI}{user interface}
\acrodef{UBC}{University of British Columbia}
\acrodef{HVAC}{heating, ventilation, and air conditioning}
\acrodef{GHG}{greenhouse gas}
\acrodef{GHGs}{greenhouse gases}
\acrodef{GIS}{geographic information systems}
\acrodef{LIDAR}{light detection and ranging}
\acrodef{UBL}{urban boundary layer}
\acrodef{UCL}{urban canopy layer}
\acrodef{RSL}{roughness sublayer}
\acrodef{LCZ}{local climate zone}
\acrodef{EC}{eddy covariance}
\acrodef{IRGA}{infrared gas analyzer}
\acrodef{PDT}{Pacific Daylight Savings Time}
\acrodef{PST}{Pacific Standard Time}
\acrodef{OSM}{OpenStreetMap}

% The acronym environment will typeset only those acronyms that were
% *actually used* in the course of the document
\begin{acronym}
\acrodef{HVAC}{heating hentilation and air conditioning}
\acrodef{GHG}{greenhouse gas}
\acrodef{GHGs}{greenhouse gases}
\acrodef{GIS}{geographic information systems}
\acrodef{LIDAR}{light detection and ranging}
\acrodef{UBL}{urban boundary layer}
\acrodef{UCL}{urban canopy layer}
\acrodef{RSL}{roughness sublayer}
\acrodef{LCZ}{local climate zone}
\acrodef{EC}{eddy covariance}
\acrodef{IRGA}{infrared gas analyzer}
\acrodef{PDT}{Pacific Daylight Savings Time}
\acrodef{PST}{Pacific Standard Time}
\acrodef{OSM}{OpenStreetMap}

\end{acronym}

% You can also use \newacro{}{} to only define acronyms
% but without explictly creating a glossary
% 
% \newacro{ANOVA}[ANOVA]{Analysis of Variance\acroextra{, a set of
%   statistical techniques to identify sources of variability between groups.}}
% \newacro{API}[API]{application programming interface}
% \newacro{GOMS}[GOMS]{Goals, Operators, Methods, and Selection\acroextra{,
%   a framework for usability analysis.}}
% \newacro{TLX}[TLX]{Task Load Index\acroextra{, an instrument for gauging
%   the subjective mental workload experienced by a human in performing
%   a task.}}
% \newacro{UI}[UI]{user interface}
% \newacro{UML}[UML]{Unified Modelling Language}
% \newacro{W3C}[W3C]{World Wide Web Consortium}
% \newacro{XML}[XML]{Extensible Markup Language}	% always input, since other macros may rely on it
	\chapter{List of Symbols and Acronyms}
% Terms:

\begin{longtable}{L{3cm}L{6cm}L{3cm}}

    % \caption[Terms]{}
    % \label{table:Terms}
% \begin{tabular}{L{3cm}L{6cm}L{3cm}} 
 \toprule 
Symbol 			&	Definition  		&	Units	\\
\midrule 
$b_a$		& 	factor for converting $\rm{m}^{-2}$ to $\rm{ha}^{-1}$ & 	i.e. $b_a = 10^{4}\,\rm{m}^{2}\,\rm{ha}^{-1}$ \\\\

$b_m$		& 	factor for converting $\rm{g}$ to $\rm{kg}$ 	& i.e. $b_m = 10^{-3}\, \rm{kg}\, \rm{g}^{-1}$ \\\\

$b_o$		& 	factor for converting $\mu\rm{mol}$ to $\rm{mol}$ & i.e. $b_m = 10^{-6}\, \mu\rm{mol}\, \rm{µmol}^{-1}$ \\\\

$b_t$ 		& 	factor for converting $\rm{s}^{-1}$ to $\rm{hr}^{-1}$ 	& i.e. $b_t = 3600\,\rm{s}\,\rm{hr}^{-1}$ \\\\

${C}$ & Rate of combustion & $\upmu\rm{mol\, m^{-2}\, s^{-1}}$ \\\\

$C_{\rm{air}}$ & heat capacity of air & $\rm{J}$ $\rm{kg}^{-1}$ $\rm{m}^{-3}$ \\\\

$c$			& 	molar concentration of $\rm{CO}_2$	&  $\upmu\rm{mol\ m^{-3}}$ \\\\

$c_{\rm{tower}}$	& $c$ at the height of the climate tower (24 m) 	& $\upmu\rm{mol\ m^{-3}}$ \\\\

$c_{\rm{mobile}}$	& $c$ at screen level height (2 m) 	& $\upmu\rm{mol\ m^{-3}}$ \\\\

EC & eddy covariance & \\\\

$F_c$ 	&	the mass flux of $\rm{CO}_2$ for a given area and time		&  $\rm{kg}\,\rm{CO}_2\,\rm{ha}^{-1}\,\rm{hr}^{-1}$ \\\\


GHG & greenhouse gas & \\\\

GIS & geographic information systems & \\\\

$H$	& sensible heat flux  	& $\rm{W} \, \rm{m}^{-2}$ \\\\

HVAC & heating hentilation and air conditioning & \\\\

IRGA & infrared gas analyzer & \\\\

LCZ & local climate zone & \\\\

LIDAR & light detection and ranging & \\\\

$M_{a}$ 	&	molecular mass of dry air		&  28.97 $\rm{g\ mol^{-1}}$ \\\\

$M_{c}$		& 	molecular mass of $\rm{CO}_2$		 & 44.01 $\rm{g\ mol^{--1}}$ \\\\

OSM & OpenStreetMap & \\\\

${P}$ & Rate of photosytnesis & $\upmu\rm{mol\, m^{-2}\, s^{-1}}$ \\\\

PDT & Pacific Daylight Savings Time \\\\
PST & Pacific Standard Time & \\\\

${R}$ & Rate of respiration & $\upmu\rm{mol\, m^{-2}\, s^{-1}}$ \\\\

$r$			&	the molar density of $\rm{CO}_2$; mole of a $\rm{CO}_2$ per mole dry air or `` molar mixing ratio'' 	&  $\upmu\rm{mol\ m^{-1}}$ or ppm \\\\

RSL & roughness sublayer& \\\\

$r_{aC}$	&  aerodynamic resistance of $\rm{CO}_2$ 	& $\rm{s} \, \rm{m}^{-1}$ \\\\

$r_{aH}$	&  aerodynamic resistance of sensible heat 	& $\rm{s} \, \rm{m}^{-1}$ \\\\

$S_c$	&	volumetric source or sink strength, describes the mass of $\rm{CO}_2$ emitted per volume and time - the $c$ that is injected or removed; normalized per ground area in this study 	&	$\upmu\rm{mol\, m^{-3}\, s^{-1}}$ \\\\

$T_{\rm{0}}$	&  surface temperature (0 m) 	& $^{\circ} \rm{C}$ \\\\

$T_{\rm{mobile}}$ &  temperature at the height of the mobile sensors (2 m) & $^{\circ} \rm{C}$ \\\\

$T_{\rm{tower}}$	&  temperature at the height of the climate tower (24 m) 	& $^{\circ}\rm{C}$ \\\\

$t$			& 	time	&  s\\\\

$u$			& 	horizontal wind speed in West-East direction & $\rm{m}\,\rm{s}^{-1}$ \\\\

UBL & urban boundary layer & \\\\

UCL & urban canopy layer & \\\\

$v$ 	& horizontal wind speed in South-North direction & $\rm{m}\,\rm{s}^{-1}$ \\\\

$w$ 	& vertical wind speed & $\rm{m}\,\rm{s}^{-1}$ \\\\

$\overline{w^{\prime}c^{\prime}}$ 	&	 molar flux 	& 	$\upmu \rm{mol}\,\rm{m}^{-2}\,\rm{s}^{-1}$\\\\

$x$			&	horizontal distance (Easting)	&  m\\\\

$y$			&	horizontal distance (Northing)	&  m\\\\

$z$			&	vertical distance (Northing)	&  m\\\\

$\rho_{a}$ 	& 	(dry) air density		&  $\rm{g\ m^{-3}}$ \\\\

$\rho\rm{CO}_2$ & mass density of $\rm{CO}_2$ &  $\rm{g}\,\rm{CO}_{2}\, \rm{m}^{-3}$\\\\

$\theta$ & potential air temperature & K \\\\

$\overline{...}$ 	&	average over a given time		&   \\\\

$\prime$	&	indicates a turblent deviation from the average	& \\\\


\bottomrule 
% \end{tabular} 
\end{longtable}















	\textspacing		% begin one-half or double spacing

	% 10. Acknowledgements (optional)
	\chapter{Acknowledgments}

Make sure to thank your friends, family, people who helped you along the way, funding agencies, etc.


	% 11. Dedication (optional)

	% Body of Thesis (not all sections may apply)
	\mainmatter

	\acresetall	% reset all acronyms used so far
	
	\chapter{hello}
	\citet{lamport-1994-ladps}
	
	\citet{bringhurst-2002-teots}
	
	\citet{Crawford:2011ir}
	
	\citet{Lietzke:2011wd}
	
	Alex wrote a paper that says he's got big muscles \citealp{Christen:2011gu}
	
	Cam wrote a paper that says Deirdre can eat 10 pies \citep{Velasco:2010tg}

	% 1. Introduction
	\chapter{Introduction}
\label{ch:Introduction}

Lorem ipsum dolor sit amet, consectetur adipiscing elit, sed do eiusmod tempor incididunt ut labore et dolore magna aliqua. Ut enim ad minim veniam, quis nostrud exercitation ullamco laboris nisi ut aliquip ex ea commodo consequat. Duis aute irure dolor in reprehenderit in voluptate velit esse cillum dolore eu fugiat nulla pariatur. Excepteur sint occaecat cupidatat non proident, sunt in culpa qui officia deserunt mollit anim id est laborum.


\begin{figure}[!htbp]
	\centering
	\includegraphics[width=0.8\textwidth]{figures/Nyan_Cat_Wide.jpg} 
	\caption[This is a nyan cat caption]{
		Insert figure caption here.
	}
	\label{fig:nyan-cat-figure}
\end{figure}\FloatBarrier

Lorem ipsum dolor sit amet, consectetur adipiscing elit, sed do eiusmod tempor incididunt ut labore et dolore magna aliqua. Ut enim ad minim veniam, quis nostrud exercitation ullamco laboris nisi ut aliquip ex ea commodo consequat. Duis aute irure dolor in reprehenderit in voluptate velit esse cillum dolore eu fugiat nulla pariatur. Excepteur sint occaecat cupidatat non proident, sunt in culpa qui officia deserunt mollit anim id est laborum.

% ------------------------------------------------------- %
\section{State of the Art}
\label{StateOfTheArt}

Lorem ipsum dolor sit amet, consectetur adipiscing elit, sed do eiusmod tempor incididunt ut labore et dolore magna aliqua. Ut enim ad minim veniam, quis nostrud exercitation ullamco laboris nisi ut aliquip ex ea commodo consequat. Duis aute irure dolor in reprehenderit in voluptate velit esse cillum dolore eu fugiat nulla pariatur. Excepteur sint occaecat cupidatat non proident, sunt in culpa qui officia deserunt mollit anim id est laborum.


% ------------------------------------------------------- %
\subsection{This is a subsection }
\label{sec:subsection1}

Lorem ipsum dolor sit amet, consectetur adipiscing elit, sed do eiusmod tempor incididunt ut labore et dolore magna aliqua. Ut enim ad minim veniam, quis nostrud exercitation ullamco laboris nisi ut aliquip ex ea commodo consequat. Duis aute irure dolor in reprehenderit in voluptate velit esse cillum dolore eu fugiat nulla pariatur. Excepteur sint occaecat cupidatat non proident, sunt in culpa qui officia deserunt mollit anim id est laborum.

% \citealp{Kellett:2013fk, Christen:2011gu}. 
% \citep{Kellett:2013fk}

\begin{table}[htbp]
\centering
    \caption[Table of measurement methods at various urban scales]{Relevant scales and monitoring methods in urban settings. Adapted from \citep{}}.
    \label{table:scales}
\begin{tabular}{L{2cm}L{1cm}L{2cm}L{3cm}L{3cm}} 
 \toprule 
Scale (Systems studies)  & Spatial Dimension & Atmospheric Layer Studied  & Temporal Resolution & Common Measurement Approaches\\
\midrule 
Micro-scale (buildings, roads, industry, greenspace) & 1-100 m  & urban canopy layer, roughness sublayer & One-time measurements at specific locations or along transects, in some cases long-term observations (5 min to years). & Traverse and vertical profile measurements in canyons, ecophysiological measurements using closed-chambers (vegetation, soils). \\
Local-scale (neighborhoods, land-use zones)          & 100 - 10 km       & internal sublayer                      & Continuous measurements that resovle diurnal and seasonal dynamics (30 min to years).                                         & Direct eddy-covariance flux measurements on towers above the city.                                                              \\
Meso-scale (cities, urban regions)                   & 10 - 100 km       & urban boundary layer                   & Short-term campaigns or continuous measurements at selected sites that resolve day-today variations and seasonal differences. & Boundary-layer budgets, upwind-downwind mixing ratio differences, regional inverse modelling, isotopic ratios.  \\            
\bottomrule 
\end{tabular} 
\end{table}
\FloatBarrier

% ------------------------------------------------------- %
\subsection{This is a subsection }
\label{sec:subsection2}


\subsubsection{This is a sub subsection}
\label{sec:subsubsection1}


\subsubsection{Objectives}
In order to address the research question, five major objectives were outlined and developed:
\begin{enumerate}
	\item Sensor Development: Develop and test a compact, mobile, and multi-modal $\rm{CO}_2$ sensor for bikes and cars.
	\item Measurement Campaign: Deploy the sensors in a targeted measurement campaign.
	\item Physical Concept: Develop a methodology to calculate emissions from measurements of $\rm{CO}_2$ mixing ratios using knowledge about atmospheric conditions. 
	\item Analysis and Evaluation: Compare the mixing ratio measurements and measured emissions to traffic and building emissions inventories. 
	% \item Visualization: Develop an interactive visualization that allows exploration of the data and concepts of the methodology and results. 
\end{enumerate}



	% 1. Methods
	\chapter{Methods}
\label{ch:Methods}


Lorem ipsum dolor sit amet, consectetur adipiscing elit, sed do eiusmod tempor incididunt ut labore et dolore magna aliqua. Ut enim ad minim veniam, quis nostrud exercitation ullamco laboris nisi ut aliquip ex ea commodo consequat. Duis aute irure dolor in reprehenderit in voluptate velit esse cillum dolore eu fugiat nulla pariatur. Excepteur sint occaecat cupidatat non proident, sunt in culpa qui officia deserunt mollit anim id est laborum.
	% results
	\chapter{Results}
\label{ch:Results}


Lorem ipsum dolor sit amet, consectetur adipiscing elit, sed do eiusmod tempor incididunt ut labore et dolore magna aliqua. Ut enim ad minim veniam, quis nostrud exercitation ullamco laboris nisi ut aliquip ex ea commodo consequat. Duis aute irure dolor in reprehenderit in voluptate velit esse cillum dolore eu fugiat nulla pariatur. Excepteur sint occaecat cupidatat non proident, sunt in culpa qui officia deserunt mollit anim id est laborum.
	% discussion
	\chapter{Conclusion}
\label{ch:Conclusion}


Lorem ipsum dolor sit amet, consectetur adipiscing elit, sed do eiusmod tempor incididunt ut labore et dolore magna aliqua. Ut enim ad minim veniam, quis nostrud exercitation ullamco laboris nisi ut aliquip ex ea commodo consequat. Duis aute irure dolor in reprehenderit in voluptate velit esse cillum dolore eu fugiat nulla pariatur. Excepteur sint occaecat cupidatat non proident, sunt in culpa qui officia deserunt mollit anim id est laborum.
	% conclusion
	\chapter{Conclusion}
\label{ch:Conclusion}


Lorem ipsum dolor sit amet, consectetur adipiscing elit, sed do eiusmod tempor incididunt ut labore et dolore magna aliqua. Ut enim ad minim veniam, quis nostrud exercitation ullamco laboris nisi ut aliquip ex ea commodo consequat. Duis aute irure dolor in reprehenderit in voluptate velit esse cillum dolore eu fugiat nulla pariatur. Excepteur sint occaecat cupidatat non proident, sunt in culpa qui officia deserunt mollit anim id est laborum.

%    5. Bibliography
\begin{singlespace}
\raggedright
\bibliographystyle{abbrvnat}
\bibliography{biblio}
\end{singlespace}

\appendix
%    6. Appendices (including copies of all required UBC Research
%       Ethics Board's Certificates of Approval)
%\include{reb-coa}	% pdfpages is useful here
\include{appendix}

\backmatter
%    7. Index
% See the makeindex package: the following page provides a quick overview
% <http://www.image.ufl.edu/help/latex/latex_indexes.shtml>


\end{document}
